Sự phát triển mạnh mẽ của công nghệ trong những năm gần đây đã mang lại nhiều lợi ích cho xã hội, đồng thời cũng tạo điều kiện cho các hình thức lừa đảo qua điện thoại ngày càng trở nên tinh vi và khó nhận biết hơn. Trong bối cảnh đó, người cao tuổi là một trong những nhóm đối tượng dễ bị tổn thương nhất do hạn chế về kỹ năng số, khả năng tiếp cận công nghệ và kinh nghiệm nhận diện các mối đe dọa trên môi trường số. Điều này đặt ra nhu cầu cấp thiết về việc xây dựng các công cụ hỗ trợ thông minh nhằm giúp người cao tuổi nâng cao khả năng tự bảo vệ trước các hành vi lừa đảo.

Từ thực tiễn trên, đề tài đã phát triển mô hình chatbot phát hiện lừa đảo điện thoại cho người cao tuổi và tiếp cận bài toán dưới góc độ Trustworthy AI nhằm đảm bảo hệ thống được phát triển theo hướng an toàn, minh bạch và có trách nhiệm đối với người sử dụng.
Đề tài đã phân tích đặc điểm của các hình thức lừa đảo điện thoại phổ biến, đồng thời xây dựng kiến trúc chatbot kết hợp giữa mô hình LLM, cơ sở tri thức (Knowledge Base) và tầng hậu kiểm HallucinationGuard. Cách tiếp cận này cho phép tận dụng khả năng hiểu ngôn ngữ tự nhiên của các mô hình AI hiện đại thông qua Groq Cloud API (Llama-3.1-8B-Instant), trong khi vẫn duy trì khả năng kiểm soát và xác thực kết quả nhằm hạn chế các rủi ro liên quan đến hallucination. Hệ thống cũng triển khai cơ chế dự phòng qua OpenRouter API và xử lý định dạng phản hồi bằng mã nguồn thay vì chỉ dựa vào hướng dẫn ban đầu.

Bên cạnh đó, đề tài đã xây dựng khung phân tích và đánh giá hệ thống dựa trên sáu trục cốt lõi của Trustworthy AI gồm Reliability, Fairness, Robustness, Explainability, Privacy và Social Impact. Đối với từng trục, các thách thức tiềm ẩn đã được phân tích một cách hệ thống, từ nguy cơ hallucination, thiên vị dữ liệu, tấn công Prompt Injection, thiếu minh bạch trong giải thích quyết định cho đến các vấn đề liên quan đến quyền riêng tư và tác động xã hội đối với nhóm người dùng yếu thế. Trên cơ sở đó, các giải pháp thiết kế phù hợp đã được đề xuất nhằm giảm thiểu rủi ro và nâng cao tính đáng tin cậy của hệ thống.

Kết quả phân tích cho thấy rằng việc phát triển một chatbot chống lừa đảo hiệu quả cho người cao tuổi không chỉ là bài toán về độ chính xác của mô hình AI mà còn là bài toán về đạo đức, trách nhiệm và sự chấp nhận của xã hội đối với công nghệ. Một hệ thống chỉ tập trung tối ưu hiệu năng kỹ thuật nhưng bỏ qua các yếu tố công bằng, minh bạch hoặc quyền riêng tư có thể tạo ra những tác động tiêu cực ngoài mong muốn. Ngược lại, việc tích hợp các nguyên tắc của Trustworthy AI ngay từ giai đoạn thiết kế giúp hệ thống trở nên an toàn hơn, dễ được chấp nhận hơn và phù hợp hơn với nhu cầu thực tế của người cao tuổi.

Mặc dù đã đạt được các mục tiêu đề ra, đề tài vẫn còn một số hạn chế. Về hạ tầng, hệ thống phụ thuộc vào Groq Cloud API với giới hạn 500.000 token/ngày, gây khó khăn cho việc kiểm thử quy mô lớn; giải pháp dự phòng OpenRouter vẫn yêu cầu kết nối Internet ổn định. Cơ sở tri thức hiện được lưu dưới dạng tệp JSON tĩnh, chưa hỗ trợ cập nhật động hoặc tích hợp nguồn cảnh báo thời gian thực, đồng thời lỗi chính tả và ký tự đặc biệt trong đầu vào vẫn có thể ảnh hưởng đến chất lượng truy xuất và suy luận. Ngoài ra, khả năng phòng thủ trước tấn công Prompt Injection còn hạn chế: hệ thống mới chỉ áp dụng cơ chế bảo vệ dựa trên hướng dẫn (instruction-based) trong prompt, chưa có tầng lọc đầu vào độc lập bằng regex, chưa tích hợp mô hình phân loại phụ trợ để phát hiện tấn công, và chưa xử lý đầu vào dưới dạng cấu trúc dữ liệu chuẩn để ngăn chặn chỉ thị độc hại lọt vào ngữ cảnh suy luận.

Về khả năng giải thích, hệ thống chưa hỗ trợ lưu vết quyết định lâu dài; log tạm thời của HallucinationGuard bị mất khi khởi động lại. Điểm tin cậy chưa được hiển thị cho người dùng, trong khi grounding score chỉ phục vụ mục đích gỡ lỗi. Khi phát hiện hallucination, hệ thống chỉ trả về thông báo chung mà chưa giải thích rõ nguyên nhân. Các dấu hiệu lừa đảo chưa có cấu trúc dữ liệu chuẩn và cũng chưa tích hợp cơ chế phản hồi từ người dùng để cải thiện mô hình.

Ngoài ra, chức năng chuyển văn bản thành giọng nói TTS chưa được triển khai do hạn chế về thời gian và kĩ thuật, dù đây là tính năng quan trọng đối với người cao tuổi có thị lực kém. Bên cạnh đó, sự thay đổi liên tục của các hình thức lừa đảo đòi hỏi hệ thống phải thường xuyên cập nhật dữ liệu và mô hình để duy trì hiệu quả phát hiện trong dài hạn.

Trong tương lai, hệ thống có thể được mở rộng theo nhiều hướng khác nhau như tích hợp nhận dạng giọng nói theo thời gian thực, bổ sung khả năng phát hiện các cuộc gọi sử dụng công nghệ giả mạo giọng nói (voice cloning), hỗ trợ đa ngôn ngữ và đa phương ngữ, đồng thời kết nối với các nguồn dữ liệu cảnh báo từ cơ quan chức năng hoặc nhà mạng. Bên cạnh đó, việc tìm hiểu sâu hơn về các phương pháp Explainable AI và các cơ chế bảo vệ quyền riêng tư tiên tiến sẽ góp phần nâng cao hơn nữa chất lượng và tính ứng dụng của giải pháp.

Tóm lại, đề tài cho thấy tiềm năng của việc ứng dụng trí tuệ nhân tạo trong hỗ trợ phòng chống lừa đảo điện thoại đối với người cao tuổi, đồng thời khẳng định tầm quan trọng của việc phát triển AI theo định hướng Trustworthy AI. Đây không chỉ là yêu cầu kỹ thuật mà còn là điều kiện cần thiết để các hệ thống trí tuệ nhân tạo có thể được triển khai một cách an toàn, hiệu quả và bền vững trong xã hội số hiện nay.
